\begingroup
\raggedright
\GaussContentsEntry{\SourceLeaf{667}\noindent\hangindent=1em Connoissance des tems ou des mouvemens célestes, à l'usage des astronomes et des navigateurs, pour l'an 1810; publiée par le bureau des longitudes. De l'imprimerie impériale 1808. Paris\GaussContentsLeaders{\mbox{.\ .\ .\ .\ .\ .\ .\ .}}\GaussContentsReference{\mbox{1809 Juni 17}\quad\mbox{Seite 515}}}
\GaussContentsEntry{\noindent\hangindent=1em \textsc{von Lindenau}: Ergebnisse von Untersuchungen über den Durchmesser der Sonne\GaussContentsLeaders{\mbox{.\ .\ .\ .\ .\ .\ .\ .\ .\ .\ .\ .\ .\ .\ .\ .\ .\ .\ .\ .\ .\ .\ .\ .}}\GaussContentsReference{\mbox{1809 Juli 17}\quad\mbox{--- 517}}}
\GaussContentsEntry{\noindent\hangindent=1em \textsc{De Rossel}: Voyage de \textsc{Dentrecasteaux}, envoyé à la recherche de la Pérouse. Publié par ordre de Sa Majesté l'Empereur et Roi. Tome second. Paris 1808\GaussContentsLeaders{\mbox{.\ .\ .\ .\ .\ .\ .\ .\ .\ .\ .\ .\ .\ .\ .\ .\ .\ .\ .\ .\ .\ .\ .}}\GaussContentsReference{\mbox{1809 Okt. 21}\quad\mbox{--- 520}}}
\GaussContentsEntry{\noindent\hangindent=1em \textsc{Regner}: Supplementum ad historiam de Parallaxeos Solaris inventione\GaussContentsLeaders{\mbox{.\ .}}\GaussContentsReference{\mbox{1810 Jan. 29}\quad\mbox{--- 523}}}
\GaussContentsEntry{\noindent\hangindent=1em \textsc{von Lindenau}: [Greenwicher Sonnenbeobachtungen von 1787---1798]\GaussContentsLeaders{\mbox{.\ .\ .}}\GaussContentsReference{\mbox{1810 Juni 16}\quad\mbox{--- 525}}}
\GaussContentsEntry{\noindent\hangindent=1em Connoissance des tems ou des mouvemens célestes à l'usage des astronomes et des navigateurs, pour l'an 1811, publiée par le bureau des longitudes. Paris 1809\GaussContentsLeaders{\mbox{.\ .\ .\ .\ .\ .\ .\ .\ .\ .\ .\ .\ .\ .\ .\ .\ .\ .\ .\ .\ .}}\GaussContentsReference{\mbox{1810 Sept. 24}\quad\mbox{--- 327}}}
\GaussContentsEntry{\noindent\hangindent=1em \textsc{Knitlmayer}: Versuch einer genauen Darstellung des Progressionsverhältnisses der Planeten- und Trabantenabstände von ihren Zentralkörpern. Brünn 1808\GaussContentsLeaders{\mbox{.\ .\ .\ .\ .\ .\ .\ .\ .\ .\ .\ .\ .\ .\ .\ .\ .\ .\ .\ .\ .\ .}}\GaussContentsReference{\mbox{1810 Nov. 10}\quad\mbox{--- 528}}}
\GaussContentsEntry{\noindent\hangindent=1em \textsc{Bessel}: Untersuchungen über die scheinbare und wahre Bahn des im Jahr 1807 erschienenen großen Kometen. Königsberg\GaussContentsLeaders{\mbox{.\ .\ .\ .\ .\ .\ .\ .\ .}}\GaussContentsReference{\mbox{1811 Feb. 16}\quad\mbox{--- 530}}}
\GaussContentsEntry{\noindent\hangindent=1em \textsc{Schulze}: Prediger in Polenz: Darstellung des Weltsystems. Ein Leitfaden für den Astronomieunterricht an Schulen. Leipzig 1811\GaussContentsLeaders{\mbox{.\ .}}\GaussContentsReference{\mbox{1811 Feb. 23}\quad\mbox{--- 533}}}
\GaussContentsEntry{\noindent\hangindent=1em \textsc{Seyffer}: Super longitudine geographica speculae astronomicae regiae, quae Monachii est, ex triginta septem defectionibus Solis observatis et ad calculos revocatis nunc primum definita. München 1810\GaussContentsLeaders{\mbox{.\ .\ .\ .\ .}}\GaussContentsReference{\mbox{1811 Mai 18}\quad\mbox{--- 533}}}
\GaussContentsEntry{\noindent\hangindent=1em \textsc{Nouet}: Description de l'Egypte. Tome premier à Paris de l'imprimerie impériale 1809. Mémoire premier. Observations astronomiques, faites en Egypte pendant les années VI, VII et VIII\GaussContentsLeaders{\mbox{.\ .\ .\ .\ .\ .\ .\ .\ .}}\GaussContentsReference{\mbox{1811 Mai 18}\quad\mbox{--- 534}}}
\GaussContentsEntry{\noindent\hangindent=1em \textsc{Pasquich}: Epitome elementorum astronomiae sphaerico-calculatoriae.\par\noindent\hangindent=1em\hspace{1em}Pars prima: Elementa theoretica astronomiae sphaerico-calculatoriae.\par\noindent\hangindent=1em\hspace{1em}Pars secunda: Elementa practica astronomiae sphaerico-calculatoriae. Wien 1811\GaussContentsLeaders{\mbox{.\ .\ .\ .\ .\ .\ .\ .\ .\ .\ .\ .\ .\ .\ .\ .\ .\ .\ .\ .\ .\ .\ .}}\GaussContentsReference{\mbox{1811 Sept. 9}\quad\mbox{--- 535}}}
\GaussContentsEntry{\noindent\hangindent=1em \textsc{Klüber}: Die Sternwarte zu Mannheim. Mannheim, Heidelberg 1811\GaussContentsLeaders{\mbox{.\ .\ .}}\GaussContentsReference{\mbox{1811 Dez. 2}\quad\mbox{--- 536}}}
\GaussContentsEntry{\noindent\hangindent=1em Connoissance des tems, etc.\ .\ .\ . pour l'an 1812. Paris 1810\GaussContentsLeaders{\mbox{.\ .\ .\ .\ .\ .}}\GaussContentsReference{\mbox{1811 Dez. 14}\quad\mbox{--- 537}}}
\GaussContentsEntry{\noindent\hangindent=1em D'Aquila: Decouverte de l'orbite de la terre du point central de l'orbite du soleil, leur situation et leur forme; de la section du zodiaque, par le plan de l'équateur, et du mouvement concordant des deux globes. Paris 1806\GaussContentsLeaders{\mbox{.\ .\ .\ .\ .\ .\ .\ .\ .\ .\ .\ .\ .\ .\ .\ .\ .\ .\ .\ .\ .\ .}}\GaussContentsReference{\mbox{1812 Juni 27}\quad\mbox{--- 541}}}
\GaussContentsEntry{\noindent\hangindent=1em \textsc{Harding}: Himmelsatlas. Lieferung III und IV.\GaussContentsLeaders{\mbox{.\ .\ .\ .\ .\ .\ .\ .\ .\ .}}\GaussContentsReference{\mbox{1812 Okt. 3}\quad\mbox{--- 541}}}
\GaussContentsEntry{\noindent\hangindent=1em \textsc{Gerling}: Methodi projectionis orthographicae usum ad calculos parallacticos facilitandos explicavit, simulque eclipsin solarem die VII. Sept. 1820 apparituram hoc modo tractatam mappaque geographica illustratam tamquam exemplum proposuit. Göttingen\GaussContentsLeaders{\mbox{.\ .\ .\ .\ .\ .\ .\ .\ .\ .\ .}}\GaussContentsReference{\mbox{1812 Nov. 28}\quad\mbox{--- 542}}}
\GaussContentsEntry{\noindent\hangindent=1em \textsc{Carlini}: Esposizione di un nuovo methodo di costruire le tavole astronomiche applicato alle tavole del sole. Mailand 1810.\GaussContentsLeaders{\mbox{.\ .\ .\ .\ .\ .\ .}}\GaussContentsReference{\mbox{1812 Dez. 3}\quad\mbox{--- 543}}}
\GaussContentsEntry{\noindent\hangindent=1em \textsc{De Zach}: Nouvelles tables d'aberration et de nutation pour quatorze cent \SourceLeaf{668}quatre étoiles, avec une table générale d'aberration pour les planètes et les comètes, précédées d'une instruction qui renferme l'explication de l'usage de ces tables, suivies de plusieurs nouvelles tables destinées à faciliter les calculs astronomiques. Marseille\GaussContentsLeaders{\mbox{.\ .\ .\ .\ .\ .\ .\ .\ .}}\GaussContentsReference{\mbox{1812 Dez. 7}\quad\mbox{Seite 546}}}
\GaussContentsEntry{\noindent\hangindent=1em \textsc{Bode}: Astronomisches Jahrbuch für das Jahr 1814. Berlin 1811\par\noindent\hangindent=1em\hspace{1em}Astronomisches Jahrbuch für das Jahr 1815. Berlin 1812\GaussContentsLeaders{\mbox{.\ .\ .\ .\ .}}\GaussContentsReference{\mbox{1813 Feb. 6}\quad\mbox{--- 548}}}
\GaussContentsEntry{\noindent\hangindent=1em Tables astronomiques publiées par le bureau des longitudes de France.\par\noindent\hangindent=1em \textsc{Burckhardt}: Tables de la lune. Paris 1812.\GaussContentsLeaders{\mbox{.\ .\ .\ .\ .\ .\ .\ .\ .\ .\ .}}\GaussContentsReference{\mbox{1813 März 8}\quad\mbox{--- 553}}}
\GaussContentsEntry{\noindent\hangindent=1em Connaissance des tems etc.\ .\ .\ . pour l'an 1813. Paris 1811.\par\noindent\hangindent=1em Connaissance des tems etc.\ .\ .\ . pour l'an 1814. Paris 1812\GaussContentsLeaders{\mbox{.\ .\ .\ .\ .\ .}}\GaussContentsReference{\mbox{1813 Juni 3}\quad\mbox{--- 554}}}
\GaussContentsEntry{\noindent\hangindent=1em \textsc{Carlini} e \textsc{Brioschi}: Effemeridi astronomische di Milano per l'anno bisestile 1812. Milano 1811.\par\noindent\hangindent=1em\hspace{1em}Effemeridi astronomiche di Milano per l'anno 1813\GaussContentsLeaders{\mbox{.\ .\ .\ .\ .\ .\ .\ .}}\GaussContentsReference{\mbox{1813 Juni 14}\quad\mbox{--- 556}}}
\GaussContentsEntry{\noindent\hangindent=1em \textsc{De Lindenau}: Investigatio nova orbitae a Mercurio circa Solem descriptae: accedunt tabulae Planetae ex elementis recens correctis et ex theoria gravitatis ill. de la Place constructae\GaussContentsLeaders{\mbox{.\ .\ .\ .\ .\ .\ .\ .\ .\ .\ .\ .}}\GaussContentsReference{\mbox{1813 Juli 5}\quad\mbox{--- 559}}}
\GaussContentsEntry{\noindent\hangindent=1em \textsc{Bessel}: Bestimmung der Präzession [Berliner Preisschrift]\GaussContentsLeaders{\mbox{.\ .\ .\ .\ .\ .}}\GaussContentsReference{\mbox{1814 Jan. 31}\quad\mbox{--- 561}}}
\GaussContentsEntry{\noindent\hangindent=1em \textsc{Pond}: Beobachtung der Sommersonnenwenden von 1812 und 1813 und der Wintersonnenwende von 1812 an der königlichen Sternwarte in Greenwich. Philosophical Transactions of the Royal Society of London for the Jear 1813\GaussContentsLeaders{\mbox{.\ .\ .\ .\ .\ .\ .\ .\ .\ .\ .\ .\ .\ .\ .\ .\ .\ .\ .\ .\ .\ .\ .}}\GaussContentsReference{\mbox{1814 Feb. 14}\quad\mbox{--- 561}}}
\GaussContentsEntry{\noindent\hangindent=1em \textsc{Bode}: Astronomisches Jahrbuch für das Jahr 1816\GaussContentsLeaders{\mbox{.\ .\ .\ .\ .\ .\ .\ .\ .\ .}}\GaussContentsReference{\mbox{1814 März 10}\quad\mbox{--- 563}}}
\GaussContentsEntry{\noindent\hangindent=1em \textsc{Carlini}: Effemeridi astronomiche di Milano per l'anno 1814\GaussContentsLeaders{\mbox{.\ .\ .\ .\ .\ .}}\GaussContentsReference{\mbox{1814 März 24}\quad\mbox{--- 566}}}
\GaussContentsEntry{\noindent\hangindent=1em \textsc{Mollweide}: Commentationes mathematico-philologicae tres, sistentes explicationem duorum locorum difficilium, alterius Virgilii, alterius Platonis itemque examinationem duorum mensurarum praeceptorum Columellae. Leipzig 1813'\GaussContentsLeaders{\mbox{.\ .\ .\ .\ .\ .\ .\ .\ .\ .\ .\ .\ .\ .\ .\ .\ .\ .}}\GaussContentsReference{\mbox{1814 Mai 2}\quad\mbox{--- 567}}}
\GaussContentsEntry{\noindent\hangindent=1em \textsc{De Zach}: L'attraction des montagnes, et ses effets sur les fils à plomb ou sur les niveaux des instrumens d'astronomie,\ .\ .\ .\ . Avignon 1814\GaussContentsLeaders{\mbox{.\ .}}\GaussContentsReference{\mbox{1814 Dez. 31}\quad\mbox{--- 569}}}
\GaussContentsEntry{\noindent\hangindent=1em \textsc{Carlini}: Effemeridi astronomiche di Milano per l'anno 1815\GaussContentsLeaders{\mbox{.\ .\ .\ .\ .\ .}}\GaussContentsReference{\mbox{1815 Jan. 19}\quad\mbox{--- 574}}}
\GaussContentsEntry{\noindent\hangindent=1em \textsc{Bode}: Astronomisches Jahrbuch für das Jahr 1817\GaussContentsLeaders{\mbox{.\ .\ .\ .\ .\ .\ .\ .\ .\ .}}\GaussContentsReference{\mbox{1815 Feb. 18}\quad\mbox{--- 576}}}
\GaussContentsEntry{\noindent\hangindent=1em Connaissance des tems etc.\ .\ .\ . pour l'an 1816\GaussContentsLeaders{\mbox{.\ .\ .\ .\ .\ .\ .\ .\ .\ .\ .}}\GaussContentsReference{\mbox{1815 März 11}\quad\mbox{--- 581}}}
\GaussContentsEntry{\noindent\hangindent=1em \textsc{Soldner}: Neue Methode zur Reduktion beobachteter Azimute\GaussContentsLeaders{\mbox{.\ .\ .\ .\ .}}\GaussContentsReference{\mbox{1815 März 23}\quad\mbox{--- 587}}}
\GaussContentsEntry{\noindent\hangindent=1em \textsc{Stefenelli}: Beiträge zur Berechnung beobachteter Azimute\GaussContentsLeaders{\mbox{.\ .\ .\ .\ .\ .}}\GaussContentsReference{\mbox{1815 März 23}\quad\mbox{--- 587}}}
\GaussContentsEntry{\noindent\hangindent=1em \textsc{Piazzi}: Praecipuarum stellarum inerrantium positiones etc. Palermo 1814\GaussContentsLeaders{\mbox{.}}\GaussContentsReference{\mbox{1815 Juni 10}\quad\mbox{--- 590}}}
\GaussContentsEntry{\noindent\hangindent=1em \textsc{Müller}: Berechnung der Sonnenfinsternisse\GaussContentsLeaders{\mbox{.\ .\ .\ .\ .\ .\ .\ .\ .\ .\ .\ .}}\GaussContentsReference{\mbox{1816 Feb. 3}\quad\mbox{--- 591}}}
\GaussContentsEntry{\noindent\hangindent=1em \textsc{Carlini}: Effemeridi astronomiche di Milano per l'anno bisestile 1816\GaussContentsLeaders{\mbox{.\ .\ .}}\GaussContentsReference{\mbox{1816 März 23}\quad\mbox{--- 592}}}
\GaussContentsEntry{\noindent\hangindent=1em \textsc{Harding}: Himmelsatlas. Lieferung V.\GaussContentsLeaders{\mbox{.\ .\ .\ .\ .\ .\ .\ .\ .\ .\ .\ .\ .\ .}}\GaussContentsReference{\mbox{1817 Jan. 9}\quad\mbox{--- 593}}}
\GaussContentsEntry{\noindent\hangindent=1em Connaissance des tems etc.\ .\ .\ . pour l'an 1818\GaussContentsLeaders{\mbox{.\ .\ .\ .\ .\ .\ .\ .\ .\ .\ .}}\GaussContentsReference{\mbox{1817 April 5}\quad\mbox{--- 595}}}
\GaussContentsEntry{\noindent\hangindent=1em \textsc{Schumacher}: De latitudine speculae Manhemiensis. Copenhagen 1816\GaussContentsLeaders{\mbox{.\ .}}\GaussContentsReference{\mbox{1817 Mai 12}\quad\mbox{--- 597}}}
\GaussContentsEntry{\noindent\hangindent=1em Tavole delle parallassi, di altezza e di latitudine calcolate dagli astronomi dell' osservatorio\ .\ .\ . nel collegio Romano. Rom 1816\GaussContentsLeaders{\mbox{.\ .\ .\ .\ .\ .\ .}}\GaussContentsReference{\mbox{1817 Juli 7}\quad\mbox{--- 599}}}
\GaussContentsEntry{\noindent\hangindent=1em \textsc{Carlini}: Effemeridi astronomiche di Milano per l'anno 1817. Mailand 1816\GaussContentsLeaders{\mbox{.}}\GaussContentsReference{\mbox{1817 Sept. 8}\quad\mbox{---\textperiodcentered{} 601}}}
\endgroup
